\chapter*{คำนำ}
\addcontentsline{toc}{chapter}{คำนำ}

ในยุคดิจิทัลที่เทคโนโลยีสารสนเทศและปัญญาประดิษฐ์ก้าวล้ำ การตรวจติดตามและส่งเสริมสุขภาวะส่วนบุคคลและองค์กร (Personal \& Organizational Health Informatics) มีบทบาทสำคัญอย่างยิ่งต่อการป้องกันโรคไม่ติดต่อเรื้อรัง (Non-Communicable Diseases หรือ NCDs) เช่น โรคเบาหวาน โรคความดันโลหิตสูง และโรคหลอดเลือดหัวใจ ซึ่งเกิดจากพฤติกรรมเนือยนิ่ง (Sedentary Behaviour) เป็นระยะเวลายาวนาน แม้ว่าสมาร์ทโฟนยุคปัจจุบันจะได้รับการติดตั้งตัวตรวจวัดความเร่งแบบสามแกน (3-Axis MEMS Accelerometer) และชิปประมวลผลการเคลื่อนไหวมาตั้งแต่โรงงาน แต่แอปพลิเคชันเพื่อสุขภาพทั่วไปมักประสบปัญหาการหยุดนับก้าวเมื่อพกสมาร์ทโฟนไว้ในกระเป๋ากางเกงหรือเมื่อปิดแอปพลิเคชันออกจากหน้าจอ

ผู้ช่วยศาสตราจารย์ ดร.ชีวะ ทัศนา ผู้ออกแบบและพัฒนา สังกัดสาขาวิชาฟิสิกส์ คณะวิทยาศาสตร์และเทคโนโลยี มหาวิทยาลัยราชภัฏรำไพพรรณี จึงได้บูรณาการองค์ความรู้ทางฟิสิกส์จลนศาสตร์ (Kinematics) การประมวลผลสัญญาณดิจิทัล (Digital Signal Processing) วิศวกรรมบริการเบื้องหลังระดับระบบปฏิบัติการ (OS Foreground Service) และปัญญาประดิษฐ์เชิงลึก เข้ากับหลักวิทยาศาสตร์การกีฬาและชีวกลศาสตร์ เพื่อออกแบบและพัฒนาชุดแอปพลิเคชันสุขภาพดิจิทัลคู่ขนานรุ่นที่ 2.0 ได้แก่ \textbf{ChewaPulse xAI (ชีวะพัลส์ เอ็กซ์เอไอ)} และ \textbf{SciRBRU Move+AI (ไซน์-อาร์บีอาร์ยู มูฟพลัสเอไอ)} ซึ่งสามารถติดตั้งใช้งานควบคู่กันบนสมาร์ทโฟนเครื่องเดียวได้โดยไม่ทับซ้อนกัน พร้อมเชื่อมต่อเข้ากับระบบฐานข้อมูลกลางเพื่อขับเคลื่อนตัวชี้วัด \textbf{OKR Key Result 3.2 (ก้าวสะสมรวมของบุคลากรในคณะวิทยาศาสตร์และเทคโนโลยี มากกว่า 10,000,000 ก้าว ในปีงบประมาณ 2570)}

นวัตกรรมสำคัญของระบบในฉบับปรับปรุงสมบูรณ์ รุ่นที่ 2.0 (v2.0) นี้ประกอบด้วย 6 มิติหลัก ได้แก่
\begin{enumerate}[label=\textbf{\arabic*.}]
  \item \textbf{มิติสถาปัตยกรรม 2 แอปพลิเคชันอิสระ (Dual Standalone Applications)} แยก Package Identifier เป็นอิสระระหว่าง \texttt{ChewaPulse xAI} (\texttt{com.rbru.health.smart\_health\_tracker}) สำหรับการดูแลสมดุลจังหวะก้าวเดินและโค้ชสุขภาพส่วนบุคคล และ \texttt{SciRBRU Move+AI} (\texttt{com.scirbru.move.ai}) สำหรับอัจฉริยภาพการเคลื่อนไหวของบุคลากรคณะวิทยาศาสตร์และเทคโนโลยี พร้อมไอคอน 3D Pixar เฉพาะแต่ละแอปพลิเคชัน
  \item \textbf{มิติการทำงานอัตโนมัติ 24 ชั่วโมงแม้ปิดแอปพลิเคชันและโหมดใส่กระเป๋ากางเกง} พัฒนาบริการเบื้องหลัง \texttt{StepForegroundService} (\texttt{START\_STICKY} ร่วมกับ \texttt{AlarmManager} และ \texttt{AutoStepBroadcastReceiver}) ที่อ่านค่าจากชิป \texttt{TYPE\_STEP\_COUNTER} คู่กับตัวกรองแรงโน้มถ่วงความถี่สูง (High-Pass Gravity Filter) สำหรับการแกว่งแบบลูกตุ้มในกระเป๋ากางเกง บันทึกข้อมูลลงหน่วยความจำถาวรแบบเรียลไทม์ทุกก้าวเดิน พร้อมแผงควบคุมให้ผู้ใช้งานกำหนดเวลาเริ่มและหยุดนับก้าวอัตโนมัติได้เอง
  \item \textbf{มิติความปลอดภัยและเสียงพูดเตือนภาษาไทยอัตโนมัติเมื่อหยุดนิ่งครบ 30 นาที} ระบบเฝ้าระวังแรงกระแทกจากการล้มฉับพลัน ($> 32.0\text{ m/s}^2$) และระบบสังเคราะห์เสียงพูดภาษาไทย (\texttt{TextToSpeech}) แจ้งเตือน \textit{``ลุกคะ ลุกค่ะ! คุณหยุดนิ่งนานเกิน 30 นาทีแล้ว''} เพื่อกระตุ้นการไหลเวียนโลหิตและป้องกันภาวะดื้ออินซูลิน
  \item \textbf{มิติระบบติดตามพิกัดดาวเทียม GPS 3 ชั้นแผนที่} รองรับการสลับชั้นแผนที่ถนนมาตรฐาน (OpenStreetMap), ภาพถ่ายดาวเทียมความละเอียดสูง (Esri World Imagery Satellite), และแผนที่ภูมิประเทศเส้นชั้นความสูง (OpenTopoMap Terrain) พร้อมคำนวณระยะทาง Haversine WGS84 ความเร็วสด อัตราเพซ และการส่งออกไฟล์ GeoJSON
  \item \textbf{มิติส่วนติดต่อผู้ใช้แบบ 4 แท็บหลัก และกราฟวิเคราะห์ย้อนหลัง 7 วัน} ออกแบบโครงสร้าง \texttt{BottomNavigationBar} 4 แท็บ (\textit{วันนี้ $\bullet$ แผนที่ดาวเทียม $\bullet$ สถิติ 7 วัน \& AI $\bullet$ โปรไฟล์}) ที่ผ่านการทดสอบความเข้ากันได้ของหน้าจอ (Multi-Viewport Responsive Tests) ตั้งแต่ขนาด $320\text{ dp}$ ถึง $1280\text{ dp}$ ไร้ปัญหาหน้าจอล้น (\texttt{RenderFlex overflow})
  \item \textbf{มิติระบบฐานข้อมูลกลาง OKR Key Result 3.2 (เป้าหมาย 10 ล้านก้าว ปีงบประมาณ 2570)} ระบบลงทะเบียนและซิงค์ก้าวเดินของบุคลากรคณะวิทยาศาสตร์และเทคโนโลยีมากกว่า 60 ท่าน (ตั้งต้น 64 ท่าน ครอบคลุมทั้ง 9 สาขาวิชา/หน่วยงาน) เข้าสู่ฐานข้อมูล SQLite3/REST API (\texttt{api.php}) แบบไม่เกิดการบวกซ้ำ (\texttt{Idempotent Upsert}) พร้อมระบบส่งออกรายงาน Excel CSV สำหรับการประเมินผลสัมฤทธิ์ระดับคณะฯ
\end{enumerate}

ผู้เขียนหวังเป็นอย่างยิ่งว่า ตำราและคู่มือปฏิบัติการฉบับสมบูรณ์เล่มนี้ จะเป็นประโยชน์ต่อคณาจารย์ บุคลากร นักศึกษา และนักพัฒนาซอฟต์แวร์ ในการนำไปประยุกต์ใช้เพื่อสร้างสรรค์นวัตกรรมสุขภาพดิจิทัลและขับเคลื่อนองค์กรสุขภาวะอย่างยั่งยืนสืบไป

\vspace{1.0cm}

\begin{flushright}
  \textbf{ผู้ช่วยศาสตราจารย์ ดร.ชีวะ ทัศนา}\\
  ผู้ออกแบบและพัฒนาระบบ ChewaPulse xAI \& SciRBRU Move+AI\\
  สาขาวิชาฟิสิกส์ คณะวิทยาศาสตร์และเทคโนโลยี\\
  มหาวิทยาลัยราชภัฏรำไพพรรณี
\end{flushright}
